\section{Binomiale}
Modella la somma di  successi in $n\in \mathbb{N}$ esperimenti di Bernoulli con identica probabilità $p$ di successo. Risulta essere quindi la somma di $n$ Bernoulliane.
\[
Bin(n,p) = \sum_1^n Be(p)
\]

\begin{figure}[H]
\centering

\begin{minipage}{0.45\textwidth}
\centering
\begin{tikzpicture}
\begin{axis}[
    axis lines = middle,
    ylabel = {$\mathbb{P}(X=k)$},
    xlabel = {$k$},
    width = 0.90\textwidth,
    height = 0.90\textwidth,
    ymin = 0, ymax = 0.22,
    xmin = -1, xmax = 26,
    legend style={at={(axis cs: 2,0.21)}, anchor=north west, font=\tiny},
    legend cell align=left,
]

% --- Bin(n=20,p=0.5) ---
\addplot[
    only marks,
    black,
    mark=*,
    mark options={fill=black, scale=0.7}
]
coordinates {
(0,0.000001)(1,0.000019)(2,0.000181)(3,0.001087)(4,0.004621)
(5,0.014786)(6,0.036965)(7,0.073929)(8,0.120134)(9,0.160179)
(10,0.176197)(11,0.160179)(12,0.120134)(13,0.073929)(14,0.036965)
(15,0.014786)(16,0.004621)(17,0.001087)(18,0.000181)(19,0.000019)
(20,0.000001)
};

% --- Bin(n=20,p=0.3) ---
\addplot[
    only marks,
    gray!60,
    mark=*,
    mark options={fill=gray!60, scale=0.7}
]
coordinates {
(0,0.000798)(1,0.006839)(2,0.027846)(3,0.071604)(4,0.130421)
(5,0.178863)(6,0.191639)(7,0.164262)(8,0.114397)(9,0.065370)
(10,0.030817)(11,0.012007)(12,0.003859)(13,0.001018)(14,0.000218)
(15,0.000037)(16,0.000005)(17,0.000001)(18,0.000000)(19,0.000000)
(20,0.000000)
};

\end{axis}
\end{tikzpicture}
\end{minipage}
\hfill
\begin{minipage}{0.45\textwidth}
\centering

\begin{tikzpicture}
\begin{axis}[
    axis lines = middle,
    ylabel = {$F(k)$},
    xlabel = {$k$},
    width = 0.90\textwidth,
    height = 0.90\textwidth,
    ymin = 0, ymax = 1.1,
    xmin = -1, xmax = 26,
    legend style={at={(axis cs: 20,0.25)}, anchor=north west, font=\tiny},
    legend cell align=left,
]

% --- CDF Bin(n=20,p=0.5) ---
\addplot[
    jump mark left,
    black,
    mark=*,
    mark options={fill=black, scale=0.7}
]
coordinates{
(0,0.000001)(1,0.000020)(2,0.000201)(3,0.001289)(4,0.005910)
(5,0.020695)(6,0.057660)(7,0.131589)(8,0.251723)(9,0.411902)
(10,0.588098)(11,0.748277)(12,0.868411)(13,0.942340)(14,0.979305)
(15,0.994090)(16,0.998711)(17,0.999799)(18,0.999980)(19,0.999999)
(20,1.000000)
};
\addlegendentry{$n=20,\, p=0.5$}

% --- CDF Bin(n=20,p=0.3) ---
\addplot[
    jump mark left,
    gray!60,
    mark=*,
    mark options={fill=gray!60, scale=0.7}
]
coordinates{
(0,0.000798)(1,0.007637)(2,0.035483)(3,0.107087)(4,0.237507)
(5,0.416370)(6,0.608009)(7,0.772271)(8,0.886668)(9,0.952038)
(10,0.982855)(11,0.994862)(12,0.998722)(13,0.999740)(14,0.999958)
(15,0.999995)(16,1.000000)(17,1.000000)(18,1.000000)(19,1.000000)
(20,1.000000)
};
\addlegendentry{$n=20,\, p=0.3$}

\end{axis}
\end{tikzpicture}

\end{minipage}

\end{figure}

\bigskip
\begin{table}[H]
\centering
\renewcommand{\arraystretch}{2}
\begin{tabular}{@{}ll@{}}
\toprule
\textbf{Supporto} 
& $\{0,1,\dots,n\}\subset\mathbb{Z}$ \\

\midrule

\textbf{Funzione di Densità} 
& $\mathbb{P}(X=k)=\binom{n}{k}p^k(1-p)^{n-k}$,\quad $k=0,\dots,n$ \\

\midrule

\textbf{Funzione di Ripartizione}
& $F(k)=\displaystyle\sum_{j=0}^{\lfloor k\rfloor}\binom{n}{j}p^j(1-p)^{n-j}$ \\

\midrule

\textbf{Valore atteso} 
& $\mathbb{E}[X]=np$ \\

\midrule

\textbf{Varianza}
& $Var(X)=np(1-p)$ \\

\midrule

\textbf{Funzione caratteristica}
& $\varphi_X(u)=\bigl[(1-p)+pe^{iu}\bigr]^n$ \\

\bottomrule
\end{tabular}
\end{table}
